%% ch02-background.tex — Chapter II: Background and Related Work
%% Shared foundation for Chapters III–VIII. Section labels sec:bg:* are
%% referenced by other chapters; do not rename them.
\chapter{Background and Related Work}\label{ch:background}


This chapter gives the shared technical foundation for the research chapters and places this dissertation in the literature. The thesis statement of \cref{ch:intro} holds that useful quantum computation in the coming decade will come from a co-designed stack. In that stack, classical high-performance computing (HPC) surrounds the quantum processing unit (QPU) at every layer. \Cref{fig:bg:stack} previews the stack and tags each layer with the research question (RQ1--RQ5) it serves. The sections follow the layers in turn:
\begin{itemize}
\item gate-model computation and the native gate sets of the processors used throughout (\cref{sec:bg:qc});
\item classical circuit simulation and the HPC systems that run it (\cref{sec:bg:sim});
\item near-term noise and its mitigation (\cref{sec:bg:nisq});
\item variational optimization and the quantum approximate optimization algorithm (\cref{sec:bg:vqa});
\item data encoding (\cref{sec:bg:encoding});
\item circuit cutting and knitting (\cref{sec:bg:cutting});
\item quantum error correction (\cref{sec:bg:qec});
\item quantum threats and the post-quantum standards that answer them (\cref{sec:bg:crypto});
\item language models for code synthesis (\cref{sec:bg:llm}).
\end{itemize}
\Cref{sec:bg:summary} closes with the gap in the literature that each research question addresses. Textbook material follows Nielsen and Chuang \cite{nielsen2010qcqi}, and the research chapters cross-reference this chapter rather than repeat it.

% alt: Block diagram of the HPC-surrounded QPU stack. A central QPU box is surrounded by eight classical-HPC boxes (problem and data, GPU simulation, LLM/RL synthesis, encoding and cutting, structured optimizers, error mitigation, QEC decoding and audit, PQC module), each tagged with the research question it serves, all enclosed in a dashed frame labelled classical HPC.
\begin{figure}[htbp]
\centering
\begin{tikzpicture}[
  font=\small,
  box/.style={draw, rounded corners=2pt, align=center, inner sep=3pt, text width=36mm, minimum height=13mm, fill=cbSky!15},
  core/.style={box, fill=cbOrange!30, very thick, minimum height=18mm},
  arr/.style={-{Stealth[length=2.2mm]}, thick, draw=cbGray},
  bi/.style={{Stealth[length=2.2mm]}-{Stealth[length=2.2mm]}, thick, draw=cbGray}]
\node[box]  (data) at (-5.1, 3.0) {Problem and data\\ {\scriptsize QUBO, images, PDE fields}};
\node[box]  (sim)  at ( 0.0, 3.0) {GPU simulation\\ {\scriptsize \cudaq{} kernels (RQ1)}};
\node[box]  (llm)  at ( 5.1, 3.0) {LLM/RL synthesis\\ {\scriptsize rubric reward (RQ4)}};
\node[box]  (enc)  at (-5.1, 0.0) {Encoding and cutting\\ {\scriptsize \qcrank{}, knitting (RQ2)}};
\node[core] (qpu)  at ( 0.0, 0.0) {\textbf{QPU}\\ {\scriptsize IBM Heron, IonQ Forte}};
\node[box]  (opt)  at ( 5.1, 0.0) {Structured optimizers\\ {\scriptsize direct angles, walks (RQ3)}};
\node[box]  (mit)  at (-5.1,-3.0) {Error mitigation\\ {\scriptsize ZNE, readout (RQ3)}};
\node[box]  (qec)  at ( 0.0,-3.0) {QEC and audit\\ {\scriptsize seeded encoders (RQ5)}};
\node[box]  (pqc)  at ( 5.1,-3.0) {PQC module\\ {\scriptsize FIPS 140-3 (RQ5)}};
\draw[arr] (data) -- (enc);
\draw[arr] (data) -- (sim);
\draw[arr] (llm) -- (sim) node[midway, above, font=\scriptsize] {circuits};
\draw[bi]  (sim) -- (qpu) node[midway, right, font=\scriptsize] {verify, calibrate};
\draw[arr] (enc) -- (qpu);
\draw[bi]  (opt) -- (qpu);
\draw[arr] (qpu) -- (qec) node[midway, right, font=\scriptsize] {syndromes};
\draw[arr] (qpu.south west) -- (mit.north east);
\draw[arr, dashed] (qec) -- (pqc) node[midway, above, font=\scriptsize] {standards};
\begin{scope}[on background layer]
\node[draw=cbGray, dashed, rounded corners=4pt, fit=(data)(llm)(mit)(pqc), inner sep=7pt] (frame) {};
\end{scope}
\node[anchor=south west, font=\footnotesize\itshape, inner sep=1pt] at (frame.north west) {Classical HPC: Perlmutter A100 GPUs, Slurm, Podman-HPC/Shifter containers};
\end{tikzpicture}
\caption{The HPC-surrounded QPU stack studied in this dissertation. Every layer around the processor is classical and runs on HPC resources; the label in parentheses gives the research question that the layer serves. Schematic.}
\label{fig:bg:stack}
\end{figure}

\section{Quantum Computation Fundamentals}\label{sec:bg:qc}

\subsection{Qubits, Gates and Measurement}

A qubit is a two-level quantum system. Its pure states are unit vectors in $\C^2$, written $\ket{\psi}=\alpha\ket{0}+\beta\ket{1}$ with $|\alpha|^2+|\beta|^2=1$. An $n$-qubit register lives in the tensor-product space $(\C^2)^{\otimes n}\cong\C^{2^n}$, and a general pure state is
\begin{equation}\label{eq:bg:statevector}
\ket{\psi}=\sum_{x\in\{0,1\}^n}\psi_x\ket{x},\qquad \sum_{x}|\psi_x|^2=1,
\end{equation}
where bit strings index the $2^n$ complex amplitudes $\psi_x$. This exponential description is the source of the hoped-for computational power and the reason that classical simulation meets a memory wall (\cref{sec:bg:sim}). Density operators $\rho$, positive semidefinite with unit trace, describe mixed states. Noise processes act as completely positive trace-preserving maps $\rho\mapsto\sum_k E_k\rho E_k^\dagger$ with Kraus operators that satisfy $\sum_k E_k^\dagger E_k=I$ \cite{nielsen2010qcqi}.

Closed-system evolution is unitary. A quantum circuit is a sequence of unitary gates, each of which acts on one or a few qubits. The gates used throughout are:
\begin{itemize}
\item the single-qubit Pauli matrices $X$, $Y$ and $Z$ and the Hadamard gate $H$;
\item the phase gates $S=\mathrm{diag}(1,i)$ and $T=\mathrm{diag}(1,e^{i\pi/4})$;
\item the rotations $R_P(\theta)=e^{-i\theta P/2}$ for $P\in\{X,Y,Z\}$;
\item the two-qubit controlled-NOT (CNOT or CX) and controlled-$Z$ (CZ) gates;
\item the parameterized interaction $e^{-i\theta Z_iZ_j/2}$, the building block of the QAOA circuits of \cref{sec:bg:vqa}.
\end{itemize}
A computation ends with a measurement in the computational basis, which returns the bit string $x$ with probability $|\psi_x|^2$ and collapses the state. One execution gives one sample. An expectation value $\expval{O}{\psi}$ of an observable $O$ must therefore be estimated from $S$ repeated executions, or shots, with a statistical error that scales as $1/\sqrt{S}$. Idealized complexity analyses hide this shot cost, a recurring theme of \cref{ch:qml,ch:synthesis}.

\subsection{Universal and Native Gate Sets}

A gate set is universal if circuits over it can approximate every unitary on $n$ qubits to arbitrary precision. The Clifford$+T$ set $\{H,S,T,\mathrm{CNOT}\}$ is universal and dominates fault-tolerant compilation (\cref{sec:bg:qec}). The Solovay--Kitaev theorem guarantees that a finite universal set approximates a single-qubit gate to accuracy $\eps$ with only $\mathrm{polylog}(1/\eps)$ gates \cite{nielsen2010qcqi}. Physical devices, however, implement a small \emph{native} gate set that their control hardware dictates. Every algorithmic circuit must be rewritten in that set before it can run.

IBM's fixed-frequency transmon processors of the Falcon and Eagle generations used the cross-resonance interaction for their entangling gate. The echoed cross-resonance (ECR) sequence cancels its dominant unwanted terms and served as the native two-qubit gate of the 127-qubit Eagle devices \cite{sheldon2016echoedcr,bravyi2022futuresc}. The Heron generation used here replaces the cross-resonance gate with a CZ gate mediated by tunable couplers. The single-qubit $\sqrt{X}$, $X$ and virtual $R_Z$ gates complete the set. Virtual $R_Z$ gates are frame changes with zero duration and zero error \cite{mckay2017virtualz}. The native set reported for the Heron devices used here is $\{\mathrm{CZ},\mathrm{ID},R_Z,\sqrt{X},X\}$, with ECR in place of CZ on the Eagle generation \cite{guo2026rubriq}.

Trapped-ion processors such as IonQ's Aria and Forte systems use the M{\o}lmer--S{\o}rensen interaction \cite{sorensen1999ms}. It couples pairs of ions through a shared motional mode and produces an $XX$-type entangling gate between arbitrary pairs. This gives all-to-all connectivity at the expense of slower gates. Native gate sets matter because the two-qubit gate count after rewriting, not the abstract count, determines the error budget of a NISQ execution.

\subsection{Transpilation and Qubit Routing}

Transpilation converts an abstract circuit into one that respects the native gate set and the coupling graph of a device. It has four parts:
\begin{enumerate}
\item gate decomposition;
\item layout, the assignment of logical qubits to physical qubits;
\item routing, the insertion of SWAP gates so that every two-qubit gate acts on physically coupled qubits;
\item gate-level optimization, such as the cancellation of adjacent inverse gates and the fusion of single-qubit rotations.
\end{enumerate}
Routing is NP-hard in general. The SABRE heuristic of Li et al. \cite{li2019sabre}, the default in \qiskit{}, performs a bidirectional search. It inserts SWAPs greedily to minimize a distance-based cost \cite{javadiabhari2024qiskit}. Each SWAP costs three CNOTs (or CZs with additional single-qubit gates) on a heavy-hex lattice. Routing can therefore double or triple the two-qubit depth of a circuit whose logical connectivity does not match the hardware. Hence the algorithm designs of \cref{ch:nisq,ch:qml} attend closely to hardware-native structure, and the hardware-aware cut selection of \cref{ch:qml} takes calibration data and physical qubit distance as first-class inputs.

\section{Classical Simulation of Quantum Circuits at HPC Scale}\label{sec:bg:sim}

\Cref{fig:bg:simlandscape} maps the simulation methods of this section, their tools and the chapters that use them.

% alt: Four columns of rounded boxes, one per simulation method: state vector, distributed state vector, matrix product state and stabilizer tableau. The top row gives the cost of each method, the middle row the tools that implement it and the bottom row the chapters that use it; an arrow underneath runs from exact general methods on the left to structure-dependent methods on the right.
\begin{figure}[htbp]
\centering
\begin{tikzpicture}[
  font=\small,
  meth/.style={draw, rounded corners=2pt, align=center, inner sep=3pt, text width=32mm, minimum height=16mm, fill=cbSky!15},
  tool/.style={draw, rounded corners=2pt, align=center, inner sep=3pt, text width=32mm, minimum height=9mm, fill=cbGreen!12},
  use/.style={draw, rounded corners=2pt, align=center, inner sep=3pt, text width=32mm, minimum height=8mm, fill=cbOrange!18},
  arr/.style={-{Stealth[length=2.2mm]}, thick, draw=cbGray}]
\node[meth] (sv)   at (-5.7, 0) {\textbf{State vector}\\ {\scriptsize $2^n$ amplitudes; exact}};
\node[meth] (dsv)  at (-1.9, 0) {\textbf{Distributed}\\ \textbf{state vector}\\ {\scriptsize sliced over GPUs}};
\node[meth] (mps)  at ( 1.9, 0) {\textbf{Matrix product}\\ \textbf{state}\\ {\scriptsize bounded entanglement}};
\node[meth] (stab) at ( 5.7, 0) {\textbf{Stabilizer tableau}\\ {\scriptsize Clifford circuits}};
\node[tool, below=3mm of sv]   (svt)   {{\scriptsize cuStateVec, \qiskit{} Aer}};
\node[tool, below=3mm of dsv]  (dsvt)  {{\scriptsize \cudaq{} multi-GPU}};
\node[tool, below=3mm of mps]  (mpst)  {{\scriptsize cuTensorNet, Aer MPS}};
\node[tool, below=3mm of stab] (stabt) {{\scriptsize Stim}};
\node[use, below=3mm of svt]   (svu)   {{\scriptsize \Cref{ch:qgear,ch:synthesis}}};
\node[use, below=3mm of dsvt]  (dsvu)  {{\scriptsize \Cref{ch:qgear}}};
\node[use, below=3mm of mpst]  (mpsu)  {{\scriptsize \Cref{ch:qml}}};
\node[use, below=3mm of stabt] (stabu) {{\scriptsize \Cref{ch:qec}}};
\draw[arr] (sv) -- (dsv);
\node[font=\scriptsize, text=cbGray, anchor=south] at ($(sv.north east)!0.5!(dsv.north west)+(0,0.6mm)$) {memory wall};
\draw[arr] ($(svu.south west)+(0,-3mm)$) -- ($(stabu.south east)+(0,-3mm)$) node[midway, below, font=\scriptsize] {exact for any circuit $\longrightarrow$ needs circuit structure};
\end{tikzpicture}
\caption{The classical simulation methods used in this dissertation: the method (top), the simulators that implement it (middle) and the chapters that use it (bottom). Schematic.}
\label{fig:bg:simlandscape}
\end{figure}

\subsection{State-Vector Simulation and the Memory Wall}

The most direct classical simulator stores all $2^n$ amplitudes of \cref{eq:bg:statevector} and applies each gate as a sparse linear map. A single-qubit gate $U=(u_{ab})$ on target qubit $t$ acts on every pair of amplitudes whose indices differ only in bit $t$,
\begin{equation}\label{eq:bg:gateapply}
\begin{pmatrix}\psi'_{x\,[t\leftarrow 0]}\\ \psi'_{x\,[t\leftarrow 1]}\end{pmatrix}
=\begin{pmatrix}u_{00}&u_{01}\\ u_{10}&u_{11}\end{pmatrix}
\begin{pmatrix}\psi_{x\,[t\leftarrow 0]}\\ \psi_{x\,[t\leftarrow 1]}\end{pmatrix}
\qquad\text{for all $x$ with bit $t$ cleared},
\end{equation}
so one gate costs $\Theta(2^n)$ arithmetic operations and one sweep over memory. A $k$-qubit gate touches $2^k$-tuples of amplitudes and costs $\Theta(2^{n+k})$. A circuit of $G$ gates therefore costs $\Theta(G\,2^n)$ time, but time is rarely the binding constraint. The state vector itself occupies
\begin{equation}\label{eq:bg:memory}
M(n)=2^n\times 16\ \text{bytes}\qquad(\text{double-precision complex amplitudes}),
\end{equation}
or half that in single precision. \Cref{fig:bg:memwall} plots \cref{eq:bg:memory} against the memory of the NVIDIA A100 accelerators \cite{choquette2021a100} used throughout. The curve doubles with every qubit: a 30-qubit state occupies \SI{17}{\giga\byte}, a 40-qubit state \SI{17.6}{\tera\byte}. The 32-qubit single-GPU, 34-qubit four-GPU and 42-qubit, 1024-GPU simulations reported for \qgear{} \cite{guo2025qgear} sit exactly where single-precision amplitudes fill the device memory. State-vector simulation at scale is thus a memory-capacity problem first and an arithmetic problem second: the memory wall behind RQ1.

% alt: Log-scale line plot of state-vector memory in gigabytes versus qubit count from 24 to 46, one solid line for double precision and one dashed line for single precision, with dotted and dash-dotted horizontal reference lines for a 40 GB A100, an 80 GB A100, a four-GPU node and 1024 GPUs, and three diamond markers at 32, 34 and 42 qubits on the single-precision line.
\begin{figure}[htbp]
\centering
\begin{tikzpicture}
\begin{axis}[
  width=0.88\linewidth, height=0.72\linewidth,
  ymode=log, xmin=24, xmax=46, ymin=0.05, ymax=2e6,
  xlabel={Number of qubits $n$}, ylabel={State-vector memory (GB)},
  xtick={24,28,...,44},
  legend pos=south east]
\addplot[cbBlue, thick, solid, mark=none] coordinates {(24,0.2684) (26,1.074) (28,4.295) (30,17.18) (32,68.72) (34,274.9) (36,1100) (38,4398) (40,1.759e+04) (42,7.037e+04) (44,2.815e+05) (46,1.126e+06)};
\addlegendentry{double precision, 16 B per amplitude}
\addplot[cbOrange, thick, dashed, mark=none] coordinates {(24,0.1342) (26,0.5369) (28,2.147) (30,8.59) (32,34.36) (34,137.4) (36,549.8) (38,2199) (40,8796) (42,3.518e+04) (44,1.407e+05) (46,5.629e+05)};
\addlegendentry{single precision, 8 B per amplitude}
\addplot[cbRed, only marks, mark=diamond*, mark size=3.2pt] coordinates {(32,34.36) (34,137.4) (42,3.518e4)};
\addlegendentry{\qgear{} runs on 1, 4 and 1024 GPUs}
\addplot[cbGray, dotted, thick, mark=none, forget plot] coordinates {(24,40) (46,40)} node[pos=1, below left, font=\scriptsize, inner sep=1pt] {one A100, \SI{40}{\giga\byte}};
\addplot[cbGray, dotted, thick, mark=none, forget plot] coordinates {(24,80) (46,80)} node[pos=0.02, above right, font=\scriptsize, inner sep=1pt] {one A100, \SI{80}{\giga\byte}};
\addplot[cbGray, dashdotted, thick, mark=none, forget plot] coordinates {(24,160) (46,160)} node[pos=0.02, above right, font=\scriptsize, inner sep=1pt] {one node, four \SI{40}{\giga\byte} GPUs};
\addplot[cbGray, dashdotted, thick, mark=none, forget plot] coordinates {(24,40960) (46,40960)} node[pos=0.02, above right, font=\scriptsize, inner sep=1pt] {1024 GPUs, \SI{40.96}{\tera\byte} in total};
\end{axis}
\end{tikzpicture}
\caption{State-vector memory footprint $2^n\times 16$ bytes (double precision) and $2^n\times 8$ bytes (single precision) as a function of the qubit count, compared with the memory of one NVIDIA A100 (40 or \SI{80}{\giga\byte}), one Perlmutter GPU node with four \SI{40}{\giga\byte} A100s, and 1024 such GPUs. Diamonds mark the 32-qubit single-GPU, 34-qubit four-GPU and 42-qubit 1024-GPU simulations reported in \cite{guo2025qgear}, placed on the single-precision curve. Schematic; reference lines are nominal device capacities.}
\label{fig:bg:memwall}
\end{figure}

\subsection{Distributed State Vectors on GPUs}

To pass the wall, the simulator must distribute the amplitudes across the memories of many devices. The standard scheme partitions on the high-order index bits. With $2^m$ devices, the $m$ most significant bits of $x$ select the device, and the remaining $n-m$ bits address the local slice. Gates on the $n-m$ low-order, or local, qubits need no communication, because both amplitudes of each pair in \cref{eq:bg:gateapply} reside on the same device. Gates on the $m$ high-order, or global, qubits pair amplitudes on different devices and need a pairwise exchange of half-slices.

A more efficient scheme permutes global and local qubits periodically through an all-to-all exchange, so that a block of subsequent gates becomes local \cite{bayraktar2023cuquantum}. On Perlmutter, the exchange crosses NVLink within a node and the HPE Slingshot network between nodes through MPI. Once the state exceeds one node, interconnect bandwidth rather than floating-point throughput bounds the achievable speed. \qgear{} \cite{guo2025qgear} relies on the multi-GPU distribution of \cudaq{} and reported close to full GPU utilization on up to 1024 A100 GPUs.

\subsection{Tensor-Network Alternatives}

Not every state requires $2^n$ numbers. A matrix product state (MPS) writes the amplitudes as a product of matrices,
\begin{equation}\label{eq:bg:mps}
\psi_{x_1x_2\cdots x_n}=\sum_{\alpha_1,\dots,\alpha_{n-1}}A^{[1]x_1}_{\alpha_1}\,A^{[2]x_2}_{\alpha_1\alpha_2}\cdots A^{[n]x_n}_{\alpha_{n-1}},
\end{equation}
where each bond index $\alpha_i$ ranges over the bond dimension $\chi$, so storage is $\bigO{n\chi^2}$ rather than $\bigO{2^n}$ \cite{perezgarcia2006mps}. The representation is exact for any state if $\chi$ may grow to $2^{\lfloor n/2\rfloor}$. It is efficient only when the entanglement across every cut of the chain is bounded, as in shallow circuits or circuits with limited long-range entanglement. Each entangling gate multiplies $\chi$ until a singular-value decomposition truncates it, which introduces a controlled approximation error. General tensor-network contraction underlies the cuTensorNet library of cuQuantum \cite{bayraktar2023cuquantum} and was the basis of the classical counter-arguments in the quantum-supremacy debate \cite{arute2019supremacy}.

\Cref{ch:qml} uses MPS compilation with a bounded bond dimension to compile sub-circuits of data-encoding circuits beyond the reach of exact simulation. The stabilizer simulators of \cref{sec:bg:qec} cover the complementary Clifford regime.

\subsection{GPU Simulators and the HPC Software Context}

\Cref{tab:bg:simulators} summarizes the simulators that recur here. \qiskit{} Aer \cite{javadiabhari2024qiskit} provides state-vector, density-matrix, MPS and stabilizer methods, with cuStateVec-backed GPU execution and batched multi-shot GPU techniques \cite{doi2023multishotgpu}. \cudaq{} \cite{kim2023cudaquantum} is NVIDIA's kernel-based programming model, and its simulation targets build on the cuQuantum SDK \cite{bayraktar2023cuquantum}: cuStateVec, with multi-GPU and multi-node distribution, and cuTensorNet.

PennyLane Lightning \cite{asadi2024pennylanelightning} is organized around adjoint differentiation for variational workloads. Qulacs \cite{suzuki2021qulacs} and QCLAB++ \cite{vanbeeumen2023qclabpp} are compact C++ simulators with single-GPU support, and Stim \cite{gidney2021stim} is a stabilizer simulator restricted to Clifford circuits. These tools expose different programming models: a \qiskit{} circuit cannot be handed to a \cudaq{} kernel without a rewrite. \Cref{ch:qgear} closes that gap without asking algorithm developers to abandon \qiskit{}.

\begin{table}[htbp]
\caption{Circuit simulators referenced in this dissertation, characterized by simulation method, parallelism and GPU support. Compiled from the cited software papers; capabilities refer to the versions described there.}
\label{tab:bg:simulators}
\centering\footnotesize
\begin{tabularx}{\linewidth}{@{}>{\raggedright\arraybackslash}p{2.6cm}>{\raggedright\arraybackslash}X>{\raggedright\arraybackslash}X>{\raggedright\arraybackslash}X>{\raggedright\arraybackslash}p{2.3cm}@{}}
\toprule
Simulator & Methods & Parallelism & GPU support & Programming model \\
\midrule
\qiskit{} Aer \cite{javadiabhari2024qiskit,doi2023multishotgpu} & state vector, density matrix, MPS, stabilizer & OpenMP threads; MPI with cuStateVec & NVIDIA via Thrust and cuStateVec; batched shots & \qiskit{} circuits (Python) \\
\addlinespace[3pt]
\cudaq{} with cuQuantum \cite{kim2023cudaquantum,bayraktar2023cuquantum} & state vector (cuStateVec), tensor network (cuTensorNet) & multi-GPU and multi-node through MPI & native NVIDIA & quantum kernels in C++ or Python \\
\addlinespace[3pt]
PennyLane Lightning \cite{asadi2024pennylanelightning} & state vector with adjoint gradients & OpenMP, Kokkos, MPI & lightning.gpu (cuStateVec), lightning.kokkos & PennyLane QNodes (Python) \\
\addlinespace[3pt]
Qulacs \cite{suzuki2021qulacs} & state vector, density matrix & OpenMP and SIMD & single GPU (CUDA) & C++ and Python \\
\addlinespace[3pt]
QCLAB++ \cite{vanbeeumen2023qclabpp} & state vector & OpenMP & single GPU (CUDA) & header-only C++ \\
\addlinespace[3pt]
Stim \cite{gidney2021stim} & stabilizer tableau, Pauli-frame sampling & SIMD bit packing & none (CPU only) & C++ and Python; Clifford circuits only \\
\bottomrule
\end{tabularx}
\end{table}

The simulations ran predominantly on Perlmutter at NERSC. Each GPU node combines an AMD EPYC processor with four NVIDIA A100 GPUs connected by NVLink. The GPUs have \SI{40}{\giga\byte} of high-bandwidth memory (\SI{80}{\giga\byte} on a subset of nodes), and the HPE Slingshot interconnect links the nodes. CPU baselines used dual-socket AMD EPYC 7763 nodes with 128 cores and \SI{512}{\giga\byte} of memory \cite{guo2025qgear,choquette2021a100}. Slurm \cite{jette2023slurm} schedules the jobs: it allocates nodes, launches MPI ranks and enforces time and memory limits, so a simulator must be packaged as a batch job.

Software environments are delivered as containers. Shifter \cite{gerhardt2017shifter} converts Docker images into a form that can be launched at scale on compute nodes. Its successor Podman-HPC \cite{stephey2022podman} provides rootless, OCI-compatible containers with GPU and MPI hooks. A \cudaq{} or PyTorch installation pinned to specific library versions can thus be rebuilt identically later, and the artifacts of \cref{app:repro} follow this practice. Mohseni et al. \cite{mohseni2024quantumsupercomputer} articulate the broader vision of a quantum supercomputer, in which QPUs form one accelerator class among many inside such a facility. That vision frames the stack of \cref{fig:bg:stack}.

\section{Noise on NISQ Hardware and Error Mitigation}\label{sec:bg:nisq}

\Cref{fig:bg:noisepipeline} places this section and \cref{sec:bg:qec} on one line, from tolerated noise to fault tolerance. This section covers the first two stages: the noise itself and its mitigation.

% alt: Flow diagram with four stages from left to right: NISQ noise, error mitigation, error correction and fault tolerance. Each stage is a rounded box that lists its main techniques, with a smaller box beneath it that names the resource the stage costs (depth budget, extra shots, extra qubits, T-count) and a label that names the chapters that work at that stage. Arrows connect the stages.
\begin{figure}[htbp]
\centering
\begin{tikzpicture}[
  font=\small,
  stage/.style={draw, rounded corners=2pt, align=center, inner sep=3pt, text width=32mm, minimum height=13mm, fill=cbSky!15},
  cost/.style={draw, rounded corners=2pt, align=center, inner sep=2pt, text width=32mm, minimum height=8mm, fill=cbOrange!18},
  arr/.style={-{Stealth[length=2.2mm]}, thick, draw=cbGray}]
\node[stage] (noise) at (-5.7, 0) {\textbf{NISQ noise}\\ {\scriptsize decoherence, gate, readout}};
\node[stage] (mit)   at (-1.9, 0) {\textbf{Error mitigation}\\ {\scriptsize ZNE, readout, PEC}};
\node[stage] (qec)   at ( 1.9, 0) {\textbf{Error correction}\\ {\scriptsize stabilizer codes}};
\node[stage] (ft)    at ( 5.7, 0) {\textbf{Fault tolerance}\\ {\scriptsize Clifford$+T$}};
\node[cost, below=3mm of noise] (noisec) {{\scriptsize Cost: depth budget}};
\node[cost, below=3mm of mit]   (mitc)   {{\scriptsize Cost: extra shots}};
\node[cost, below=3mm of qec]   (qecc)   {{\scriptsize Cost: extra qubits}};
\node[cost, below=3mm of ft]    (ftc)    {{\scriptsize Cost: $T$-count}};
\draw[arr] (noise) -- (mit) node[midway, yshift=8mm, font=\scriptsize, text=cbGray] {tolerate};
\draw[arr] (mit) -- (qec) node[midway, yshift=8mm, font=\scriptsize, text=cbGray] {mitigate};
\draw[arr] (qec) -- (ft) node[midway, yshift=8mm, font=\scriptsize, text=cbGray] {correct};
\node[below=2mm of noisec, font=\scriptsize] {\Cref{tab:bg:devices}};
\node[below=2mm of mitc, font=\scriptsize] {\Cref{ch:nisq}};
\node[below=2mm of qecc, font=\scriptsize] {\Cref{ch:qec}};
\node[below=2mm of ftc, font=\scriptsize] {\Cref{ch:synthesis}};
\end{tikzpicture}
\caption{From tolerated noise to fault tolerance: each stage, the resource it costs, and the chapter that works at that stage. Schematic.}
\label{fig:bg:noisepipeline}
\end{figure}

\subsection{Noise Channels and Device Metrics}

Preskill's term noisy intermediate-scale quantum (NISQ) \cite{preskill2018nisq} designates processors with tens to a few hundred qubits whose errors are not corrected but must be tolerated. Two incoherent processes dominate. Energy relaxation has the time $T_1$, over which the excited-state population decays as $e^{-t/T_1}$. Dephasing has the time $T_2\le 2T_1$, over which the off-diagonal elements of the density matrix decay as $e^{-t/T_2}$. Idle qubits also decohere: a qubit that waits for a time $t$ while others are operated on accumulates an error of order $t/T_1$ and $t/T_2$. This idling decoherence is why long, sparse circuits are not free and why delay-aware scheduling matters.

Randomized benchmarking \cite{magesan2012rb} quantifies gate errors. On current superconducting devices, single-qubit errors are of order $10^{-4}$ and two-qubit errors of order $10^{-3}$. An example is the median CZ error of \num{1.46e-3} on ibm\_pittsburgh during the experiments of \cref{ch:qml} \cite{guo2025shardq}. Readout error is the probability that a qubit prepared in $\ket{0}$ or $\ket{1}$ is reported as the other value. It is typically of order $10^{-2}$, an order of magnitude worse than gate error.

Crosstalk denotes unwanted coupling between simultaneously driven or neighbouring qubits, including residual $ZZ$ interactions with spectator qubits. Isolated gate benchmarks do not capture it, and it degrades circuits whose gates run in parallel \cite{sheldon2016echoedcr}. \Cref{tab:bg:devices} lists the processors used here with the values that the source publications report.

\begin{table}[htbp]
\caption{Quantum processors used in this dissertation. Qubit counts and two-qubit error figures are those reported in the publications of the chapters that used each device (a median unless marked as a mean or a bound: \texttt{ibm\_kingston} from \cite{guo2025vqt}, \texttt{ibm\_pittsburgh} from \cite{guo2025shardq}, \texttt{ibm\_boston} from \cite{guo2026nnqa}); native gate sets follow \cite{guo2026rubriq}. Cells not reported in those sources are left blank.}
\label{tab:bg:devices}
\centering\footnotesize
\begin{tabular}{@{}llccll@{}}
\toprule
Device & Family & Qubits & Native 2q gate & Median 2q error & Used in \\
\midrule
ibm\_torino & IBM Heron r1 & 133 & CZ & --- & \cref{ch:nisq} \\
ibm\_marrakesh & IBM Heron r2 & 156 & CZ & --- & \cref{ch:nisq,ch:qml} \\
ibm\_kingston & IBM Heron r2 & 156 & CZ & $\approx\num{2.5e-3}$ (mean) & \cref{ch:qml} \\
ibm\_pittsburgh & IBM Heron r3 & 156 & CZ & \num{1.46e-3} & \cref{ch:nisq,ch:qml} \\
ibm\_boston & IBM Heron r3 & 156 & CZ & $<\num{5e-3}$ & \cref{ch:qml,ch:synthesis} \\
ibm\_miami & IBM Nighthawk & 120 & CZ & --- & \cref{ch:qml} \\
IonQ Aria-1 & trapped ion & 25 & MS & --- & \cref{ch:qml} \\
IonQ Forte-1 & trapped ion & 36 & MS & --- & \cref{ch:qml,ch:synthesis} \\
\bottomrule
\end{tabular}
\end{table}

Consider a circuit with $G_2$ two-qubit gates of error $\eps_2$, $G_1$ single-qubit gates of error $\eps_1$ and $n$ measured qubits with readout error $\eps_r$. A simple heuristic puts its success probability at roughly $(1-\eps_2)^{G_2}(1-\eps_1)^{G_1}(1-\eps_r)^{n}$ before idling and crosstalk are counted. With $\eps_2\approx10^{-3}$, the two-qubit budget alone falls below $e^{-1}$ at about a thousand gates.

In practice, coherent errors and crosstalk make the usable depth considerably smaller. The \deal{} experiments of \cref{ch:nisq} observed performance degrading beyond roughly 41--43 CZ gates on Heron processors \cite{guo2025deal}. Depth is the enemy on NISQ hardware, and every design in this dissertation is a strategy for buying accuracy with less depth.

\subsection{Error Mitigation}

Error mitigation reduces the bias of expectation values estimated on noisy hardware at the cost of additional circuit executions, without the qubit overhead of error correction. Zero-noise extrapolation (ZNE) was introduced independently by Temme et al. \cite{temme2017errormitigation} and Li and Benjamin \cite{li2017activeerror}. It executes the circuit at several amplified noise levels $\lambda_1<\lambda_2<\dots<\lambda_m$, with $\lambda=1$ the native level, and extrapolates the measured expectation values $E(\lambda_i)$ to $\lambda=0$. Under the model $E(\lambda)=E^{\ast}+\sum_{k=1}^{m-1}a_k\lambda^k$, Richardson extrapolation gives
\begin{equation}\label{eq:bg:zne}
E^{\ast}\approx\sum_{i=1}^{m}\gamma_iE(\lambda_i),\qquad \sum_{i}\gamma_i=1,\qquad \sum_{i}\gamma_i\lambda_i^{k}=0\quad(k=1,\dots,m-1),
\end{equation}
which for two points reduces to $E^{\ast}\approx(\lambda_2E(\lambda_1)-\lambda_1E(\lambda_2))/(\lambda_2-\lambda_1)$.

Kandala et al. \cite{kandala2019errormitigation} amplified noise by stretching gate pulses, which requires pulse-level access. Digital ZNE \cite{giurgicatiron2020digitalzne} amplifies noise at the gate level instead. Unitary folding replaces a circuit $U$, or a subset of its gates, by $U(U^{\dagger}U)^{m}$, which is logically the identity but multiplies the gate noise by an odd factor $\lambda=2m+1$. Identity or delay insertion lengthens idle periods, so the coherence-limited part of the noise is scaled separately from the gate-limited part. The adaptive digital ZNE of \cref{ch:nisq} \cite{guo2025deal} belongs to the latter family. It inserts delay gates whose durations come from device calibration data, so crosstalk and coherence errors are controlled rather than merely amplified.

The variance of the extrapolated estimator grows with $\sum_i\gamma_i^2$, so ZNE trades bias for shots. ZNE also fails when the noise is not a smooth function of $\lambda$, which is one reason the extrapolation must be validated against simulation. Measurement errors are mitigated separately. The readout channel is modelled as a stochastic matrix $A$ with $A_{yx}=\Pr[\text{read }y\mid\text{prepared }x]$, calibrated from preparation circuits. Then $\hat{p}_{\mathrm{ideal}}=A^{-1}\hat{p}_{\mathrm{noisy}}$, or a constrained least-squares variant, recovers the ideal distribution. Because $A$ has $2^n\times2^n$ entries, scalable methods assume a tensor-product or few-qubit-correlated structure \cite{bravyi2021measurement}.

Probabilistic error cancellation \cite{temme2017errormitigation} inverts the gate noise channel as a quasi-probability distribution and is unbiased. Its sampling overhead, however, grows exponentially with the number of noisy gates, so ZNE and readout mitigation are the practical choice at the depths considered here.

\section{Variational Algorithms and QAOA}\label{sec:bg:vqa}

\subsection{Variational Quantum Algorithms}

Variational quantum algorithms (VQAs) \cite{cerezo2021vqa} pair a parameterized quantum circuit $U(\bm{\theta})$ with a classical optimizer. The QPU prepares $\ket{\psi(\bm{\theta})}=U(\bm{\theta})\ket{0}^{\otimes n}$ and estimates a cost $C(\bm{\theta})=\expval{H}{\psi(\bm{\theta})}$ from shots. The classical side updates $\bm{\theta}$, and the loop repeats. The circuits are shallow by design, which suits NISQ hardware. The variational quantum eigensolver targets ground-state energies of molecular or spin Hamiltonians. QAOA targets combinatorial problems, and quantum neural networks target supervised learning.

All share the same failure modes: the number of cost evaluations that the optimizer demands, each of which costs shots and queue time, and the flat optimization landscape discussed below. Many evaluations use slightly different parameters, so GPU simulation is the natural development environment. This dissertation uses it to validate variational circuits before hardware submission and, in \cref{ch:synthesis}, as the reward oracle of a learning loop.

\subsection{QUBO Problems and the QAOA Ansatz}

Quadratic unconstrained binary optimization (QUBO) asks for $\bm{x}^{\ast}=\argmin_{\bm{x}\in\{0,1\}^n}\bm{x}^{\mathsf{T}}Q\bm{x}$ for a symmetric matrix $Q\in\R^{n\times n}$. Glover et al. \cite{glover2018qubo} survey how penalty terms fold constraints into $Q$. Monta{\~n}ez-Barrera et al. \cite{montanezbarrera2024unbalanced} show how inequality constraints can be encoded without slack variables. Typical examples are:
\begin{itemize}
\item MaxCut on a graph $(V,E)$, the canonical example, with $Q$ built from the adjacency matrix;
\item the travelling-salesman problem, which uses $n^2$ binary variables $x_{v,t}$ (city $v$ at position $t$) with one-hot penalties on rows and columns;
\item the knapsack problem, which penalizes capacity violations;
\item portfolio optimization, which minimizes the risk $\bm{x}^{\mathsf{T}}\Sigma\bm{x}$ minus the expected return under a budget penalty.
\end{itemize}
The substitution $x_i=(1-Z_i)/2$ converts the QUBO objective into a diagonal cost Hamiltonian
\begin{equation}\label{eq:bg:hcost}
\Hcost=\sum_{i<j}J_{ij}Z_iZ_j+\sum_{i}h_iZ_i+c,\qquad J_{ij}=\tfrac12Q_{ij},\quad h_i=-\tfrac12\sum_{j}Q_{ij},
\end{equation}
whose ground state encodes $\bm{x}^{\ast}$. The quantum approximate optimization algorithm (QAOA) \cite{farhi2014qaoa} prepares the trial state
\begin{equation}\label{eq:bg:qaoa}
\ket{\bm{\gamma},\bm{\beta}}=\prod_{k=1}^{p}e^{-i\beta_k\Hmix}\,e^{-i\gamma_k\Hcost}\ket{+}^{\otimes n},\qquad \Hmix=\sum_{i=1}^{n}X_i,
\end{equation}
with $2p$ angles, and minimizes $\expval{\Hcost}{\bm{\gamma},\bm{\beta}}$. Each cost layer is a product of commuting $ZZ$ and $Z$ rotations, and every $e^{-i\gamma Z_iZ_j}$ compiles to two CNOTs and one $R_Z(2\gamma)$. Each mixer layer is a product of $R_X(2\beta)$ gates, as \cref{fig:bg:qaoa} shows. As $p\to\infty$, the ansatz can reproduce adiabatic evolution and reach the exact optimum. On hardware, the two-qubit gate count $2p|E|$ plus routing overhead binds long before that.

Weidenfeller et al. \cite{weidenfeller2022qaoascaling} found that SWAP overhead dominates QAOA on heavy-hex hardware for dense problem graphs. Sachdeva et al. \cite{sachdeva2024qaoa127} reported that a 127-qubit gate-model device could outperform annealers on structured instances. Abbas et al. \cite{abbas2024challengesqopt} survey the open problems of the field critically. Continuous-time quantum walks \cite{childs2009quantumwalk} offer an alternative view of the mixer as a walk on the hypercube and inspire the approximate-walk formulation of \cref{ch:nisq}.

% alt: Three-qubit quantum circuit showing one QAOA layer: Hadamard gates on all qubits, a ZZ rotation on qubits 0 and 1 decomposed into CNOT, RZ of two gamma and CNOT, a boxed ZZ rotation on qubits 1 and 2, then RX of two beta mixers on every qubit followed by measurements.
\begin{figure}[htbp]
\centering
\begin{quantikz}[row sep={0.75cm,between origins}, column sep=0.3cm]
\lstick{$q_0=\ket{0}$} & \gate{H} & \ctrl{1}\gategroup[3,steps=4,style={dashed,rounded corners,inner xsep=2pt},background,label style={label position=below,anchor=north,yshift=-0.2cm}]{cost layer $e^{-i\gamma_k\Hcost}$} & & \ctrl{1} & & \gate{R_X(2\beta_k)}\gategroup[3,steps=1,style={dashed,rounded corners,inner xsep=2pt},background,label style={label position=below,anchor=north,yshift=-0.2cm}]{mixer} & \meter{} \\
\lstick{$q_1=\ket{0}$} & \gate{H} & \targ{} & \gate{R_Z(2\gamma_k)} & \targ{} & \gate[2]{e^{-i\gamma_k Z_1Z_2}} & \gate{R_X(2\beta_k)} & \meter{} \\
\lstick{$q_2=\ket{0}$} & \gate{H} & & & & & \gate{R_X(2\beta_k)} & \meter{}
\end{quantikz}
\caption{One QAOA layer for a three-qubit cost Hamiltonian with couplings on qubit pairs $(0,1)$ and $(1,2)$. The first $ZZ$ rotation is shown in its native decomposition into two CNOTs and one $R_Z(2\gamma_k)$; the second is drawn as a block. Single-qubit $Z_i$ terms of \cref{eq:bg:hcost} would add $R_Z$ gates to the cost layer. Schematic.}
\label{fig:bg:qaoa}
\end{figure}

\subsection{Trainability and Parameter Initialization}

The classical optimization of $\bm{\gamma}$ and $\bm{\beta}$ is the Achilles' heel of QAOA. McClean et al. \cite{mcclean2018barren} studied sufficiently random parameterized circuits, which approximate unitary 2-designs. For them, the variance of every gradient component vanishes exponentially in $n$, $\mathrm{Var}[\partial_{\theta}C]\in\bigO{b^{-n}}$ for some $b>1$. Gradients are then exponentially small almost everywhere: the barren plateau.

Holmes et al. \cite{holmes2022expressibility} connected the severity of the plateau to ansatz expressibility. Wang et al. \cite{wang2021noisebp} showed that local noise induces a plateau whose severity grows with circuit depth, independently of the ansatz. This ties trainability directly to the depth budget of \cref{sec:bg:nisq}. Each gradient component also costs shots, and optimizer iterations multiply the hardware queue time.

The parameter-initialization literature responds by exploiting structure. Zhou et al. \cite{zhou2020qaoa} observed that optimal QAOA angles concentrate and vary smoothly with $p$, which permits interpolation-based initialization from smaller $p$. Egger et al. \cite{egger2021warmstart} warm-start QAOA from the continuous relaxation of the QUBO, which biases the initial state and mixer toward a classical solution. Fixed-angle and parameter-transfer approaches reuse angles across instances of the same problem class. All of these still treat the angles as free parameters to be tuned.

The \deal{} approach of \cref{ch:nisq} \cite{guo2025deal} goes one step further: it maps the QUBO coefficients directly onto ansatz angles, so a deterministic construction replaces optimization. The \qawa{} method \cite{guo2025qawa} replaces the black-box loop with a shallow, classically traceable oracle whose mid-circuit measurement statistics are interpretable without full state tomography.

\section{Encoding Classical Data into Quantum States}\label{sec:bg:encoding}

\subsection{Basis, Angle and Amplitude Encodings}

Any quantum algorithm that processes classical data must first load it into a quantum state. The choice of encoding determines both the circuit cost and what can be computed afterwards \cite{schuld2021dataencoding,larose2020encodings}. Basis encoding writes a bit string into the computational basis, $\bm{b}\mapsto\ket{\bm{b}}$. It is exact and cheap but uses one qubit per bit. Angle encoding maps each real feature $x_j$ to a rotation of one qubit, $\ket{\psi(\bm{x})}=\bigotimes_j R_Y(x_j)\ket{0}$, so $n$ features cost $n$ qubits and constant depth. Repeated or higher-frequency angle encodings act as feature maps whose Fourier spectrum controls the expressivity of the resulting model \cite{schuld2021dataencoding,havlicek2019featurespaces}.

Amplitude encoding stores $N=2^n$ normalized values as the amplitudes $\psi_x$ of an $n$-qubit state and achieves exponential compression in qubits. But preparing an arbitrary amplitude-encoded state requires $\Theta(2^n)$ gates in general, and reading the values back requires tomography. The compression is paid for at both ends. Proposals that assume amplitude-encoded inputs at no cost, including the quantum linear-systems algorithm \cite{harrow2009hhl}, are exposed to this input problem \cite{biamonte2017qml}.

\subsection{QPIXL and QCrank}

Image and array data motivated a family of encodings. They keep the exponential compression of amplitude encoding for the \emph{address} but store the \emph{values} as rotation angles. The QPIXL framework of Amankwah et al. \cite{amankwah2022qpixl} shows that the flexible representation of quantum images and related schemes reduce to a uniformly controlled rotation. After a Walsh--Hadamard transform of the angles and Gray-code ordering of the controls, its circuit uses only $R_Y$ and CNOT gates. It admits compression when small transformed angles are discarded.

\qcrank{} \cite{balewski2024qcrank} generalizes this to $n_d$ data qubits controlled by $n_a$ address qubits. It encodes a sequence of $N=2^{n_a}n_d$ real values $\theta_{a,d}$ in two steps: it prepares the address register in uniform superposition and applies to each data qubit $d$ the uniformly controlled rotation
\begin{equation}\label{eq:bg:ucr}
U_d=\sum_{a=0}^{2^{n_a}-1}\ket{a}\!\bra{a}\otimes R_Y(\theta_{a,d}),\qquad
\ket{\Psi}=\frac{1}{\sqrt{2^{n_a}}}\sum_{a}\ket{a}\bigotimes_{d=1}^{n_d}\Bigl(\cos\tfrac{\theta_{a,d}}{2}\ket{0}+\sin\tfrac{\theta_{a,d}}{2}\ket{1}\Bigr).
\end{equation}
The M{\"o}tt{\"o}nen decomposition \cite{mottonen2004ucr} implements each $U_d$ with exactly $2^{n_a}$ single-qubit $R_Y$ rotations interleaved with $2^{n_a}$ CNOTs whose controls follow a Gray code, as \cref{fig:bg:qcrank} sketches. The transformed angles $\tilde{\bm{\theta}}=M^{-1}\bm{\theta}$ come from a Walsh--Hadamard-type matrix $M$. All $n_d$ data qubits are processed in parallel, so the total CNOT count is $\Theta(N)$ and the circuit depth is $\Theta(2^{n_a})$, independent of $n_d$.

% alt: QCrank circuit with two address qubits and one data qubit. Hadamard gates act on the address qubits; the data qubit receives an alternating sequence of four RY rotations with transformed angles and four CNOTs controlled by the address qubits in Gray-code order; all three qubits are measured at the end.
\begin{figure}[htbp]
\centering
\begin{quantikz}[row sep={0.75cm,between origins}, column sep=0.28cm]
\lstick{$a_0=\ket{0}$} & \gate{H} & & & \ctrl{2} & & & & \ctrl{2} & \meter{} \\
\lstick{$a_1=\ket{0}$} & \gate{H} & \ctrl{1} & & & & \ctrl{1} & & & \meter{} \\
\lstick{$d=\ket{0}$} & \gate{R_Y(\tilde{\theta}_0)} & \targ{} & \gate{R_Y(\tilde{\theta}_1)} & \targ{} & \gate{R_Y(\tilde{\theta}_2)} & \targ{} & \gate{R_Y(\tilde{\theta}_3)} & \targ{} & \meter{}
\end{quantikz}
\caption{\qcrank{} encoding of $N=4$ real values on $n_a=2$ address qubits and $n_d=1$ data qubit. The uniformly controlled rotation of \cref{eq:bg:ucr} is decomposed following \cite{mottonen2004ucr} into $2^{n_a}=4$ rotations by the transformed angles $\tilde{\theta}_j$ and four CNOTs whose controls follow a Gray code; additional data qubits repeat the same pattern in parallel. Schematic after \cite{balewski2024qcrank}.}
\label{fig:bg:qcrank}
\end{figure}

Readout is equally parallel. Measuring all qubits and conditioning on the address bits yields $\expval{Z_d}_a=\cos\theta_{a,d}$. Every stored value is thus recovered from the same pool of shots as an expectation value rather than by tomography. With $S$ total shots, each address is sampled about $S/2^{n_a}$ times, and the standard error of each recovered value scales as $\sqrt{2^{n_a}/S}$. This shot cost of reading out is the price of the parallelism, and the cost model of \cref{ch:synthesis} quantifies it \cite{guo2026nonlinearwaves}.

Balewski et al. demonstrated \qcrank{} on trapped-ion and transmon hardware \cite{balewski2024qcrank} and compiled it for a dynamically programmable neutral-atom processor \cite{balewski2025qcrankdpqa}. \qcrank{} circuits are a benchmark of \cref{ch:qgear}, the substrate of the attention and encoding primitives of \cref{ch:qml}, and the objects that \shardq{} cuts and reassembles.

\subsection{Polynomial Computation, Block Encodings and Readout Cost}

Once real values are encoded as angles, arithmetic can be performed on them coherently. The EHands protocol of Balewski et al. \cite{balewski2025ehands} provides primitives for multiplication, weighted addition and polynomial evaluation on real-valued encoded states. A polynomial of the encoded value is thus produced without intermediate measurement. This is the starting point for the neural-native arithmetic of \cref{ch:qml} \cite{guo2026nnqa}, which compiles a trained polynomial network into such primitives. Block encoding is a different route to matrix arithmetic. A unitary $U_A$ on $a+n$ qubits block-encodes an $n$-qubit operator $A$ with normalization $\alpha$ if $(\bra{0}^{\otimes a}\otimes I)\,U_A\,(\ket{0}^{\otimes a}\otimes I)=A/\alpha$.

The FABLE compiler \cite{camps2022fable} produces approximate block encodings of dense matrices with the same uniformly controlled rotation structure as \qcrank{}. Quantum singular value transformation (QSVT) \cite{gilyen2019qsvt} then applies a polynomial transformation to the singular values of $A$. It needs a number of block-encoding calls equal to the polynomial degree, and it unifies Hamiltonian simulation, linear-system solving and amplitude amplification. The vectorized dot product of \cref{ch:qml} \cite{guo2025vqt} can be read as an approximate block encoding whose output is recovered as expectation values.

In all these schemes, the honest cost includes the shots needed to read the result. An expectation value estimated to precision $\eps$ requires $\bigO{1/\eps^2}$ shots. When all $N$ outputs must be read, as in the nonlinear time stepping of \cref{ch:synthesis}, the readout cost scales as $\bigO{N/\eps^2}$ and can cancel the coherent speed-up \cite{guo2026nonlinearwaves}. Classical shadows \cite{huang2020shadows} reduce the cost of estimating many observables from randomized measurements. They are the natural comparison point for such parallel readout.

\section{Circuit Cutting and Knitting}\label{sec:bg:cutting}

Circuit cutting partitions a large circuit into sub-circuits small enough for the available hardware or simulator. The sub-circuits are executed separately, and classical post-processing reconstructs the output of the original circuit, a step often called knitting. Two kinds of cut exist.

A \emph{wire cut} severs a qubit wire between two gates. Peng et al. \cite{peng2020circuitcutting} showed that the identity channel on the cut wire can be expanded in a basis of measure-and-prepare operations. These measure in the Pauli bases and prepare Pauli eigenstates. The original expectation value then becomes a signed sum of expectation values of sub-circuits.

A \emph{gate cut} removes a two-qubit gate. Mitarai and Fujii \cite{mitarai2021virtualgate} showed that a CNOT or CZ can be written as a quasi-probability decomposition (QPD)
\begin{equation}\label{eq:bg:qpd}
\mathcal{U}=\sum_{i}c_i\,\mathcal{F}_i,\qquad c_i\in\R,\qquad \gamma=\sum_{i}|c_i|,
\end{equation}
into local single-qubit operations $\mathcal{F}_i$ with real coefficients that may be negative. Sampling the terms with probability $|c_i|/\gamma$ and weighting the outcomes by $\gamma\,\mathrm{sign}(c_i)$ yields an unbiased estimator whose variance is inflated by $\gamma^2$ relative to the uncut circuit. For a CNOT $\gamma=3$, so cutting $k$ gates independently incurs a sampling overhead of $\gamma^{2k}=3^{2k}$, and a wire cut incurs $4^2=16$ per cut. This exponential overhead is the fundamental limitation of the technique and the reason that cut placement is an optimization problem in its own right.

Bravyi, Smith and Smolin \cite{bravyi2016trading} analyzed the general trade between classical and quantum resources in such decompositions. Piveteau and Sutter \cite{piveteau2024knitting} showed that classical communication between the sub-circuits reduces the overhead of cutting $n$ parallel CNOTs from $\gamma=3^{n}$ to $\gamma=2^{n+1}-1$. They also showed that gate cuts with communication are strictly more efficient than wire cuts.

CutQC \cite{tang2021cutqc} was the first end-to-end system. It finds cuts with a mixed-integer program, runs the sub-circuits on small devices and reconstructs full probability distributions with dynamic-definition post-processing. It demonstrated the evaluation of circuits larger than the device. IBM's qiskit-addon-cutting packages QPD-based gate and wire cutting for production use. The \shardq{} framework of \cref{ch:qml} \cite{guo2025shardq} specializes cutting to \qcrank{}-type data-encoding circuits. Its SparseCut heuristic chooses gate cuts from physical qubit distance and calibrated error rates. Its reconstruction is global rather than pairwise, with the $\bigO{3^{2k}}$ overhead absorbed on GPUs.

\section{Quantum Error Correction and Fault Tolerance}\label{sec:bg:qec}

\subsection{Stabilizer Codes}

Quantum error correction (QEC) encodes logical information redundantly across many physical qubits, so that errors can be detected and undone without a measurement of the logical state. The stabilizer formalism of Gottesman \cite{gottesman1997stabilizer} describes nearly all codes in practical use. Let $\mathcal{P}_n$ be the $n$-qubit Pauli group and $\mathcal{S}\subset\mathcal{P}_n$ an abelian subgroup not containing $-I$, generated by $n-k$ independent elements $S_1,\dots,S_{n-k}$. The code space is the joint $+1$ eigenspace
\begin{equation}\label{eq:bg:stabilizer}
\mathcal{C}=\bigl\{\ket{\psi}\in(\C^2)^{\otimes n}:\ S\ket{\psi}=\ket{\psi}\ \ \forall S\in\mathcal{S}\bigr\},\qquad \dim\mathcal{C}=2^{k},
\end{equation}
which encodes $k$ logical qubits. The logical operators are the elements of the normalizer $N(\mathcal{S})$ that are not in $\mathcal{S}$, and the code distance $d$ is the minimum weight of such an element. The code is written $\code{n}{k}{d}$ and corrects arbitrary errors on any $t=\lfloor(d-1)/2\rfloor$ qubits.

A Pauli error $E$ that anticommutes with some generator flips the measured eigenvalue of that generator. The vector of eigenvalues, the syndrome, identifies $E$ up to stabilizers. The general condition for a set of errors $\{E_a\}$ to be correctable is the Knill--Laflamme condition $PE_a^{\dagger}E_bP=c_{ab}P$ on the code projector $P$ \cite{knill1997theory}.

The smallest code that corrects an arbitrary single-qubit error is the $\code{5}{1}{3}$ perfect code of Laflamme et al. \cite{laflamme1996perfect}. It saturates the quantum Hamming bound: its $2^4=16$ syndromes exactly label the $3\times5=15$ weight-one Pauli errors plus the trivial one. \Cref{ch:qec} exploits this property and enumerates all $3n$ weight-one faults exhaustively. Steane's $\code{7}{1}{3}$ code \cite{steane1996codes} is a Calderbank--Shor--Steane (CSS) code \cite{calderbank1996goodcodes} built from the classical $[7,4,3]$ Hamming code, with separate $X$-type and $Z$-type generators. The CSS structure gives it transversal Clifford gates. Terhal \cite{terhal2015qecmemories} reviews the theory of QEC for quantum memories.

\subsection{Surface Codes, Bicycle Codes and Thresholds}

Kitaev's toric code \cite{kitaev2003anyons} and its planar variant, the surface code \cite{dennis2002topological,fowler2012surfacecodes}, place qubits on a two-dimensional lattice. Weight-four $X$-type and $Z$-type checks sit on alternating faces, so every check involves only nearest neighbours. Syndrome extraction therefore needs only local CNOTs and one ancilla per check. A distance-$d$ patch uses roughly $2d^2$ qubits. The threshold theorem states that below a physical error rate $p_{\mathrm{th}}$, an increase of $d$ suppresses the logical error rate as $p_L\approx A\,(p/p_{\mathrm{th}})^{\lceil d/2\rceil}$. For the surface code under circuit-level depolarizing noise, $p_{\mathrm{th}}$ is of order $1\%$ \cite{fowler2012surfacecodes}.

Decoding maps a syndrome history to a likely error, usually by minimum-weight perfect matching. Sundaresan et al. \cite{sundaresan2023subsystem} compared matching and maximum-likelihood decoders on a heavy-hex subsystem code over multiple rounds. The 2025 experiment of Google Quantum AI \cite{acharya2025belowthreshold} operated distance-3, -5 and -7 surface codes below threshold. Each increase of the distance by two reduced the logical error rate by a factor $\Lambda\approx2$.

The rate $k/n$ of the surface code vanishes as $d$ grows, so quantum low-density parity-check codes with higher rate are attractive. The bivariate bicycle codes of Bravyi et al. \cite{bravyi2024bicycle}, exemplified by the $\code{144}{12}{12}$ code with weight-six checks, are one family. They promise roughly a tenfold reduction in overhead relative to the surface code at comparable noise, at the cost of non-planar connectivity.

Syndrome extraction circuits are themselves faulty, and a fault on an ancilla can propagate to several data qubits. Fault-tolerant extraction schedules and repeated measurement rounds are designed to contain such faults. Instead of correcting, many small experiments post-select: runs with a nontrivial syndrome are rejected. This lowers the error rate of the accepted runs at the cost of discarded shots. The security analysis of \cref{ch:qec} \cite{guo2026qecaudit} is built precisely around what survives such rejection.

\subsection{Clifford Simulation, Random Ensembles and Magic}

The Clifford group $\mathcal{C}_n$ is the normalizer of the Pauli group, and $H$, $S$ and CNOT generate it. The Gottesman--Knill theorem states that circuits of Clifford gates, Pauli measurements and stabilizer-state inputs can be simulated efficiently on a classical computer. The $n$ stabilizer generators, rather than $2^n$ amplitudes, track the state. Aaronson and Gottesman \cite{aaronson2004stabilizer} gave a tableau representation with $\bigO{n}$ cost per gate and $\bigO{n^2}$ per measurement.

Stim \cite{gidney2021stim} implements this with bit-packed, vectorized tableaus and a fast Pauli-frame sampler for repeated noisy runs. It samples syndrome-extraction circuits with thousands of qubits and is the standard tool for estimating logical error rates and thresholds. The gate-level simulations of \cref{ch:qec} at $n=17$ use it.

Encoding circuits of stabilizer codes are Clifford, which has two consequences exploited here. First, ensembles of random Clifford unitaries are cheap to sample and to simulate, whereas Haar-random unitaries on $n$ qubits require exponentially many gates. Second, the Clifford group is an exact unitary 2-design \cite{dankert2009designs}, and in fact a 3-design \cite{webb2016clifford3design}: for every operator $X$ on two copies of the system,
\begin{equation}\label{eq:bg:design}
\mathop{\mathbb{E}}_{U\sim\mathcal{C}_n}\bigl[U^{\otimes2}X(U^{\dagger})^{\otimes2}\bigr]=\mathop{\mathbb{E}}_{U\sim\mathrm{Haar}}\bigl[U^{\otimes2}X(U^{\dagger})^{\otimes2}\bigr],
\end{equation}
so any quantity that depends on at most the second moment of the ensemble takes its Haar value when averaged over Cliffords. This includes average fidelities and the disturbance measures of \cref{ch:qec}. Classical shadows \cite{huang2020shadows} rely on the same fact.

Clifford gates alone are not universal. The $T$ gate supplies the missing ingredient, often called magic. In the surface code, $T$ cannot be implemented transversally. It is instead injected through magic-state distillation \cite{bravyi2005magic}, whose cost dominates fault-tolerant resource estimates. The $T$-count is therefore the primary cost driver of a fault-tolerant circuit, ahead of Clifford count and depth. This is why it is a weighted term in the synthesis rubric of \cref{ch:synthesis} \cite{guo2026rubriq}, and why $T$-consuming arithmetic dominates the cryptanalytic estimates of \cref{sec:bg:crypto}.

\section{Quantum Threats, Post-Quantum Cryptography and Standards}\label{sec:bg:crypto}

\subsection{Cryptanalytic Quantum Algorithms and Resource Estimates}

Shor's algorithm \cite{shor1997factoring} factors integers and computes discrete logarithms in polynomial time. It uses the quantum Fourier transform to find the period of modular exponentiation. This breaks RSA, finite-field Diffie--Hellman and elliptic-curve cryptography. Grover's algorithm \cite{grover1996search} finds a marked item among $N$ with $\bigO{\sqrt{N}}$ queries and thereby halves the effective key length of symmetric ciphers and hash functions. The standard response is to double key sizes rather than to change algorithms. Resource estimates set the gap between these asymptotic statements and an actual attack.

Gidney and Eker{\aa} \cite{gidney2021rsa} estimated that factoring a 2048-bit RSA modulus would take about eight hours on 20 million noisy superconducting qubits. Their estimate assumed a surface code at a physical error rate of $10^{-3}$, and the $T$-gate-consuming modular arithmetic dominated. Gidney's 2025 update \cite{gidney2025million} lowers the estimate to fewer than one million noisy qubits running for less than a week. It keeps the same assumptions: a $0.1\%$ gate error rate, a square grid with nearest-neighbour connections, a \SI{1}{\micro\second} surface-code cycle and a \SI{10}{\micro\second} reaction time. The reduction comes from approximate residue arithmetic, yoked surface codes for idle logical qubits and magic-state cultivation in place of distillation. These estimates connect the fault-tolerant cost accounting of \cref{sec:bg:qec} to concrete policy timelines.

Mosca's inequality \cite{mosca2018cybersecurity} states the following. Let $x$ be the time for which data must remain secret, $y$ the time needed to migrate infrastructure and $z$ the time until a cryptographically relevant quantum computer exists. If $x$ plus $y$ exceeds $z$, the data is already at risk, because an adversary can harvest encrypted traffic now and decrypt it later.

\subsection{Standardization and Validation}

Post-quantum cryptography (PQC) \cite{bernstein2017pqc} replaces number-theoretic assumptions with problems believed to be hard for quantum computers, chiefly structured lattices, hash functions and error-correcting codes. NIST opened its standardization process with a public call for algorithms in December 2016 \cite{nist2016pqccall}. After three rounds of public cryptanalysis, it announced the selected algorithms in July 2022 \cite{nist2022ir8413}. In August 2024, it published three Federal Information Processing Standards:
\begin{itemize}
\item FIPS 203, the module-lattice key-encapsulation mechanism ML-KEM derived from CRYSTALS-Kyber \cite{nist2024fips203};
\item FIPS 204, the module-lattice digital signature ML-DSA derived from CRYSTALS-Dilithium \cite{nist2024fips204};
\item FIPS 205, the stateless hash-based signature SLH-DSA derived from SPHINCS$^{+}$ \cite{nist2024fips205}.
\end{itemize}
A fourth signature standard based on FALCON and an additional code-based key-encapsulation mechanism, HQC, are in progress.

Standardizing an algorithm is not the same as trusting an implementation. FIPS 140-3 \cite{nist2019fips1403}, aligned with ISO/IEC 19790, defines four security levels for cryptographic modules and the Cryptographic Module Validation Program. Under that program, an accredited laboratory tests the approved algorithms, key management, self-tests and physical security of a module. Only then may the module be used in federal systems. Among the approved components, SP 800-90A \cite{nist2015sp80090a} specifies deterministic random bit generators (DRBGs): Hash\_DRBG, HMAC\_DRBG and CTR\_DRBG with their instantiation, reseeding and health-test requirements.

Randomness quality is a validation item in its own right, because every PQC key and signature depends on it. The Light Rider QPC SDK of \cref{ch:pqc} \cite{guo2026lightriderqpc} is a FIPS 140-3-style PQC module. The randomness audit of \cref{ch:qec} \cite{guo2026qecaudit} deliberately borrows the seeding and reseeding vocabulary of SP 800-90A.

\subsection{Threat Models for Cloud Quantum Hardware}

Quantum key distribution (QKD) addresses the key-exchange problem physically rather than computationally. BB84 \cite{bennett1984bb84} encodes key bits in non-orthogonal photon states, so eavesdropping disturbs them detectably. Device-independent QKD \cite{vazirani2014diqkd} certifies security from Bell-inequality violations without trusting the devices. Wiretap-channel analyses such as \cite{pan2020secretkey} bound the secret-key rate under restricted eavesdropping. QKD secures links rather than computations and is complementary to PQC.

A newer question is the security of computation on shared, cloud-accessed QPUs, where a user's circuit runs on hardware the user neither owns nor observes. Blind quantum computation \cite{broadbent2009blind} hides the computation from the server. Mahadev's protocol \cite{mahadev2018verification} lets a classical client verify the output of a quantum server under cryptographic assumptions. Neither is practical on current devices.

At the hardware level, multi-programming, or co-tenancy, executes circuits of different users simultaneously on one chip. Ash-Saki et al. \cite{ashsaki2020crosstalk} showed that a co-tenant can exploit crosstalk to degrade or bias a victim's circuit. This established fault injection through shared physical resources as a threat vector. The audit of \cref{ch:qec} \cite{guo2026qecaudit} formalizes such an adversary under a bounded fault model. It asks how the public, fixed structure of QEC encoders exposes them. This is an emerging topic at the intersection of error correction and systems security.

\section{Large Language Models for Code and Circuit Synthesis}\label{sec:bg:llm}

Large language models (LLMs) trained on source code generate syntactically valid programs from natural-language or partial-code prompts. Codex \cite{chen2021codex} established the functional-correctness evaluation methodology of the field: generated programs are executed against held-out tests rather than compared textually to a reference. Open models such as Qwen2.5-Coder \cite{hui2024qwen25coder} now reach strong performance at 7--32 billion parameters. Quantum circuits expressed in \qiskit{}, OpenQASM or \cudaq{} are programs, so the same models can be prompted to synthesize circuits.

Circuit synthesis is in one respect an unusually favourable target for such models, because a classical simulator can score a generated circuit exactly against its specification. The fidelity to a target state or unitary, adherence to a coupling graph and gate counts such as the $T$-count of \cref{sec:bg:qec} are all computable, so the simulator doubles as an oracle. Supervised fine-tuning alone, however, does not teach a model to respect such quantitative constraints. Generated code that fails to parse or to compile is common.

Reinforcement learning from verifiable rewards supplies the missing feedback. Proximal policy optimization \cite{schulman2017ppo} maximizes a clipped surrogate objective and requires a learned value network. Group relative policy optimization (GRPO), introduced with DeepSeekMath \cite{shao2024deepseekmath}, removes the value network. It samples a group of $N$ completions per prompt and uses the group-normalized reward $\hat{A}_i=(r_i-\bar{r})/\sigma_r$ as the advantage. DeepSeek-R1 \cite{deepseek2025r1} showed that GRPO with rule-based rewards can induce long-form reasoning.

Two techniques make the fine-tuning of a multi-billion-parameter model affordable. Low-rank adaptation (LoRA) \cite{hu2022lora} freezes the pretrained weights $W_0$ and learns a low-rank update $W=W_0+BA$ with $\mathrm{rank}(BA)=r\ll\min(d_{\mathrm{in}},d_{\mathrm{out}})$. This reduces the number of trainable parameters by orders of magnitude. DeepSpeed ZeRO \cite{rajbhandari2020zero} partitions optimizer states (stage~1), gradients (stage~2) and parameters (stage~3) across data-parallel GPUs. The memory footprint of a training run thus shrinks with the number of devices. Generative alternatives to LLMs exist: F{\"u}rrutter et al. \cite{furrutter2024genqc} synthesize circuits with diffusion models conditioned on the target unitary.

The \rubriq{} system of \cref{ch:synthesis} \cite{guo2026rubriq} combines these ingredients: a 7-billion-parameter coder with rank-64 LoRA, GRPO with a group size of eight and ZeRO stage 2 across A100 nodes. A programmatic rubric reward, evaluated by \cudaq{} simulation inside the loop, places the GPU simulator of \cref{sec:bg:sim} at the centre of the learning process.

\section{Positioning of This Dissertation}\label{sec:bg:summary}

The preceding sections identify, for each layer of the stack in \cref{fig:bg:stack}, a body of established results and an unresolved gap. We close with those gaps stated in terms of the five research questions of \cref{ch:intro}.

\textbf{RQ1 (Scale).} GPU state-vector simulators and distributed frameworks exist (\cref{sec:bg:sim}), and \cudaq{} in particular reaches thousand-GPU scale, but each requires programs written in its own model. \qiskit{}-native workflows, which dominate algorithm development and hardware access, either run on Aer with limited multi-node scaling or must be ported by hand. No published pipeline takes \qiskit{} circuits unchanged to distributed GPU simulation with containerized, Slurm-ready deployment. None characterizes where communication overtakes computation for typical algorithm circuits at 30--42 qubits on a production HPC system. \Cref{ch:qgear} addresses this gap with \qgear{} \cite{guo2025qgear,guo2025qgearsoftware}.

\textbf{RQ2 (Data).} The encodings of \cref{sec:bg:encoding} make loading and parallel readout efficient. Those encodings and learnable hybrid models have rarely been combined. \qcrank{}-style vectorized encodings had not been made learnable or used as attention primitives. Nor had they been cut hardware-aware into sub-circuits or used as the target of an exact compilation of a trained network. Gradient-based hybrid models, conversely, were evaluated without accounting for shots or for the QPU-versus-GPU gradient trade-off. \Cref{ch:qml} addresses this gap with \qpie{} \cite{guo2025qpie}, \vqt{} \cite{guo2025vqt}, \shardq{} \cite{guo2025shardq} and \nnqa{} \cite{guo2026nnqa}.

\textbf{RQ3 (Noise).} The depth budget of \cref{sec:bg:nisq} and the optimization cost and barren plateaus of \cref{sec:bg:vqa} limit QAOA on hardware. Warm starts and parameter concentration reduce but do not eliminate the black-box loop. ZNE is usually applied as a generic post-processing step disconnected from the ansatz. Missing is a formulation in which the problem coefficients determine the angles directly, the ansatz is built from the native entangling gate of the device, and noise amplification is designed jointly with the circuit. Also missing is an oracle whose outputs are interpretable without tomography. \Cref{ch:nisq} addresses this gap with \deal{} \cite{guo2025deal} and \qawa{} \cite{guo2025qawa}.

\textbf{RQ4 (Synthesis).} LLM- and RL-based program synthesis (\cref{sec:bg:llm}) had not been coupled to the fault-tolerant cost metrics of \cref{sec:bg:qec} through a simulator-in-the-loop reward. Diffusion-based circuit synthesis targets unitaries rather than constraint-aware programs. At the same time, quantum algorithms for nonlinear partial differential equations hide the measurement-and-reload cost of \cref{sec:bg:encoding} inside linear embeddings \cite{jin2024schrodingerization}. No end-to-end account of when a quantum solver pays off exists. \Cref{ch:synthesis} addresses this gap with \rubriq{} \cite{guo2026rubriq,guo2026rubriqsoftware} and with the split-step cost model for nonlinear waves \cite{guo2026nonlinearwaves,guo2026splitstepqude}.

\textbf{RQ5 (Trust).} QEC theory (\cref{sec:bg:qec}) assumes an adversary-free noise model. The cloud threat models of \cref{sec:bg:crypto} have been studied for NISQ circuits but not for encoders and syndrome-based rejection. Nor had QEC been connected to the randomness and module-validation standards that govern classical cryptography. \Cref{ch:qec} addresses this gap with the structured-randomness audit \cite{guo2026qecaudit}. \Cref{ch:pqc} addresses it with a FIPS 140-3-style post-quantum module \cite{guo2026lightriderqpc} positioned against the ML-KEM, ML-DSA and SLH-DSA standards \cite{nist2024fips203,nist2024fips204,nist2024fips205}, FIPS 140-3 \cite{nist2019fips1403} and SP 800-90A \cite{nist2015sp80090a}.

Together, these five gaps trace the outline of \cref{fig:bg:stack}. A QPU is useful only to the extent that the classical layers around it are fast, structured, data-efficient, cost-aware and trustworthy. The chapters that follow build those layers in turn.
